\section*{Capítulo \ref{ch:b2:precambrian-earth} --- La Tierra precámbrica}

\begin{solution}{exo:b2:precambrian-earth:1}
(a) \omterm{def:b2:precambrian-earth:eons}{Hádico} $4567 - 4031 = \qty{536}{Myr}$; \omterm{def:b2:precambrian-earth:eons}{Arcaico} $4031 - 2500 = \qty{1531}{Myr}$;
\omterm{def:b2:precambrian-earth:eons}{Proterozoico} $2500 - 538.8 = \qty{1961}{Myr}$. (b) $(4567 - 538.8)/4567 = 0.882$: unos siete
octavos. (c) $4567 - 4404 = \qty{163}{Myr}$ para el circón; $4567 - 4030 = \qty{537}{Myr}$ para
las rocas más antiguas.
\end{solution}

\begin{solution}{exo:b2:precambrian-earth:2}
(a) Los factores $e^{\lambda t}$ con $\lambda = \ln2/T$ y $t = \qty{3}{Gyr}$ son 1.593
(\ce{^{238}U}), 19.18 (\ce{^{235}U}), 1.160 (\ce{^{232}Th}) y 5.292 (\ce{^{40}K}): las
contribuciones eran 3.053, 1.573, 2.424 y \qty{4.329}{pW}, \qty{11.38}{pW} en total, frente a
\qty{4.906}{pW} hoy. (b) $11.38/4.906 = 2.32$. (c) Hoy domina el torio 232; entonces dominaba el
potasio 40.
\end{solution}

\begin{solution}{exo:b2:precambrian-earth:3}
(a) $1616 + 18.0\times(28 - 32.2) = \qty{1540}{\celsius}$. (b) $32.2 + (1500 - 1616)/18.0 =
\qty{25.8}{\percent}$. (c) Un análisis de roca total no es la composición de un líquido: una
colada de \omterm{def:b2:precambrian-earth:greenstone}{komatiita} puede contener cristales de olivino acumulados a partir del líquido, que añaden MgO. Su
\qty{38}{\percent} queda más allá del \qty{32}{\percent} de los líquidos komatiíticos más ricos, y
más allá de los experimentos en los que se apoya la recta.
\end{solution}

\begin{solution}{exo:b2:precambrian-earth:4}
Rocas más antiguas, \qty{4.03}{Ga}; primeros estromatolitos, \qty{3.43}{Ga}; final de los domos de Pilbara,
\qty{3.1}{Ga}; \omterm{def:b2:precambrian-earth:goe}{Gran Evento de Oxidación}, de unos 2.4 a \qty{2.1}{Ga}; últimas \omterm{def:b2:precambrian-earth:greenstone}{komatiitas} comunes,
\qty{1.9}{Ga}; segundo aumento del oxígeno, hacia \qty{0.8}{Ga}; final de la marinoense,
\qty{0.635}{Ga}; primeros fósiles ediacáricos, hacia \qty{0.57}{Ga}.
\end{solution}

\begin{solution}{exo:b2:precambrian-earth:5}
(a) $\tfrac14\times1361\times0.706 = \qty{240.2}{W.m^{-2}}$ absorbidos y $\sigma\times288^4 =
\qty{390.1}{W.m^{-2}}$, así que $\varepsilon = 240.2/390.1 = 0.616$. (b) $\tfrac14\times1361\times0.4
= \qty{136.1}{W.m^{-2}}$, $T^4 = 136.1/(0.616\times\num{5.670e-8})$, $T = \qty{249.8}{K}$.
(c) A \qty{275.5}{K} el albedo es $0.6 - 0.306\times12.5/20 = 0.4088$: $R = 340.25\times0.5912 =
\qty{201.2}{W.m^{-2}}$ y $E = 0.616\,\sigma\times275.5^4 = \qty{201.2}{W.m^{-2}}$. (d) $R' =
340.25\times0.306/20 = \qty{5.21}{W.m^{-2}.K^{-1}}$, $E' = 4\varepsilon\sigma T^3 =
\qty{2.92}{W.m^{-2}.K^{-1}}$: $R' > E'$, así que el equilibrio es inestable
(\cref{prop:b2:precambrian-earth:ebm}).
\end{solution}

\begin{solution}{exo:b2:precambrian-earth:6}
(a) $1/(1 + 0.4\times2.4/4.567) = 0.826$; $1/(1 + 0.4\times0.65/4.567) = 0.946$. (b) El estado
cálido necesita al menos 0.932 de la luz solar actual: imposible hace \qty{2.4}{Ga}, posible
hace \qty{0.65}{Ga}. (c) $0.826/0.946 = 0.873$: un \qty{13}{\percent} más débil, cerca del
\qty{10}{\percent} citado, que procede de otro modelo solar.
\end{solution}

\begin{solution}{exo:b2:precambrian-earth:7}
(a) Esturtiense $716.9 - 658.9 = \qty{58.0}{Myr}$; marinoense $644.5 - 635.4 = \qty{9.1}{Myr}$;
interglaciar $658.9 - 644.5 = \qty{14.4}{Myr}$. (b) $(58.0 + 9.1)/85 = 0.79$: unas cuatro quintas partes
del periodo. (c) Su comienzo solo está acotado entre 649.9 y \qty{639.0}{Ma}: con esas
cotas, la marinoense duró entre unos 3 y \qty{15}{Myr}.
\end{solution}

\begin{solution}{exo:b2:precambrian-earth:8}
(a) $\tau = \frac{0.05}{0.95}\times\frac{44.01}{28.97}\times\frac{\num{5.1e18}}{\num{0.26e12}} =
0.05263\times1.519\times\num{1.962e7} = \qty{1.57}{Myr}$. (b) $0.05/\num{278e-6} = 180$ veces.
(c) El doble, \qty{3.14}{Myr}.
\end{solution}

\begin{solution}{exo:b2:precambrian-earth:9}
(a) Los cantos cayeron en un barro blando y curvaron sus láminas hacia abajo: son cantos caídos, liberados por
hielo flotante. (b) Un hielo que llega al mar a \ang{10} del ecuador es la huella esperada de una
bola de nieve, y muy difícil de obtener si solo están heladas las regiones polares. (c) Un \omterm{def:b2:precambrian-earth:snowball}{carbonato de cobertera}:
registra el final repentino de la glaciación, con una meteorización intensa y la inundación de las
plataformas por el agua de fusión.
\end{solution}

\begin{solution}{exo:b2:precambrian-earth:10}
(a) $0.94 > 0.932$, el punto de inflexión cálido: el equilibrio cálido, $288\times0.94^{1/4} =
\qty{283.6}{K}$, queda justo por encima de \qty{283}{K}; y $0.94 < 1.227$, así que también existe la bola de nieve.
(b) $249.8\times0.94^{1/4} = \qty{246.0}{K}$. (c) $\varepsilon = 0.94\times1361\times0.4/(4\sigma
\times263^4) = 511.7/1085.1 = 0.472$. (d) $\sigma\times263^4 = \qty{271.3}{W.m^{-2}}$, y $(0.616
- 0.472)\times271.3 = \qty{39}{W.m^{-2}}$.
\end{solution}

\begin{solution}{exo:b2:precambrian-earth:11}
(a) $S_{\text{m}}/S_{\text{f}} = (T_1/T_2)^4\,(1 - a_{\text{w}})/(1 - a_{\text{i}})$: la
emisividad se cancela. $(263/283)^4\times0.706/0.4 = 0.7459\times1.765 = 1.317$. (b) Un contraste
de albedo menor entre el hielo y el agua abierta, como un hielo más oscuro y polvoriento, o una rampa de
temperatura más ancha. (c) Una bola de nieve refleja tanta luz solar que sigue fría con cualquier luz solar o
efecto invernadero con los que una Tierra cálida se helaría; para deshelarla, el forzamiento debe superar al de la
congelación en toda la anchura del bucle.
\end{solution}

\begin{solution}{exo:b2:precambrian-earth:12}
Cerca de $T^*$, $C\dot\delta = -(E' - R')\delta$, así que $\delta \propto e^{-t/\tau}$ con $\tau =
C/(E' - R')$. En el equilibrio cálido $R' = 0$ y $E' = 4\times0.616\times\num{5.670e-8}\times
288^3 = \qty{3.34}{W.m^{-2}.K^{-1}}$; $C = 1025\times3992\times50 = \qty{2.05e8}{J.m^{-2}.K^{-1}}$,
así que $\tau = \qty{6.1e7}{s}$, unos \qty{1.9}{yr}. En este modelo el equilibrio cálido mantiene $R' =
0$ hasta que alcanza \qty{283}{K}, de modo que $\tau$ sigue cerca de dos años hasta el punto de inflexión; después
el estado desaparece. Con un albedo que varía de forma continua, $R'$ crece a medida que el equilibrio se acerca al
pliegue, donde $E' = R'$: $\tau$ crece sin límite, una ralentización que anuncia el punto de
inflexión.
\end{solution}

\begin{solution}{pb:b2:precambrian-earth:1}
\textbf{1.} $\num{0.0203e-6}\times0.9929\times\num{95.13e-6} = \qty{1.92}{pW.kg^{-1}}$
(\ce{^{238}U}); $\num{0.0203e-6}\times0.00711\times\num{568.47e-6} = \qty{0.0820}{pW.kg^{-1}}$
(\ce{^{235}U}); $\num{0.0795e-6}\times\num{26.28e-6} = \qty{2.09}{pW.kg^{-1}}$ (\ce{^{232}Th});
$\num{240e-6}\times\num{1.20e-4}\times\num{28.47e-6} = \qty{0.820}{pW.kg^{-1}}$ (\ce{^{40}K}).
\textbf{2.} \qty{4.91}{pW.kg^{-1}}.
\textbf{3.} $\ln2/T$: 0.155, 0.985, 0.0495 y \qty{0.555}{Gyr^{-1}}.
\textbf{4.} Factores 1.52, 14.3, 1.14 y 4.48: 2.91, 1.17, 2.39 y \qty{3.67}{pW.kg^{-1}}.
\textbf{5.} \qty{10.1}{pW.kg^{-1}}, 2.07 veces el actual.
\textbf{6.} El potasio 40. Su periodo de semidesintegración, \qty{1.25}{Gyr}, es corto: desde entonces han pasado más de dos periodos,
y se ha desintegrado más de tres cuartas partes del \ce{^{40}K} de aquella época.
\textbf{7.} $1/(1 + 0.4\times2.7/4.567) = 0.809$; $0.809\times1361 = \qty{1101}{W.m^{-2}}$.
\textbf{8.} $288\times0.809^{1/4} = \qty{273}{K}$, por debajo de \qty{283}{K}: no es un estado cálido válido;
la Tierra se helaría.
\textbf{9.} $(1101\times0.4/(4\times0.616\,\sigma))^{1/4} = \qty{237}{K}$.
\textbf{10.} $0.616\times0.809 = 0.498$.
\textbf{11.} De $0.004/\num{278e-6} = 14$ a $0.03/\num{278e-6} = 108$ veces.
\textbf{12.} El termostato de la meteorización de los silicatos: la meteorización se ralentiza en un clima frío, lo que deja
acumularse el dióxido de carbono volcánico, y se acelera en uno cálido.
\textbf{13.} 0.826 de la actual, por debajo de 0.932: no hay estado cálido; la Tierra con el aire actual sería una
bola de nieve.
\textbf{14.} $\frac{0.1}{0.9}\times1.519\times\num{1.962e7} = \qty{3.31}{Myr}$.
\textbf{15.} $\frac{0.2}{0.8}\times1.519\times\num{1.962e7} = \qty{7.45}{Myr}$.
\textbf{16.} El \omterm{def:b2:precambrian-earth:snowball}{carbonato de cobertera}, una dolomía laminada que se apoya bruscamente sobre los depósitos glaciares,
precipitado en un mar alimentado por la meteorización intensa de las secuelas cálidas e inundado por el
agua de fusión.
\textbf{17.} El hierro(II) solo permanece disuelto en agua sin oxígeno: el hierro se acumula en el
océano anóxico bajo el hielo y precipita como hierro(III) allí donde llega el oxígeno, por ejemplo con el
agua de fusión.
\textbf{18.} $(1616 - 1450)/(32.2 - 23) = \qty{18.0}{\celsius}$ por punto porcentual.
\textbf{19.} $1616 + 18.0\times(30 - 32.2) = \qty{1576}{\celsius}$.
\textbf{20.} \qty{1616}{\celsius}, el punto medido.
\textbf{21.} $1616 - 18.0\times20.2 = \qty{1252}{\celsius}$.
\textbf{22.} $1616 - 20\times20.2 = \qty{1212}{\celsius}$: \qty{40}{\celsius} menos, una décima parte de
la diferencia entre las lavas; la conclusión no depende de la elección.
\textbf{23.} Las \omterm{def:b2:precambrian-earth:greenstone}{komatiitas} procedían de plumas, más calientes que el manto ambiente, y de grados de fusión
grandes en profundidad; la temperatura de una lava no es la temperatura potencial de su fuente. El
manto en su conjunto era de 200 a \qty{300}{K} más caliente.
\textbf{24.} La lava komatiítica más caliente estaba unos \textbf{\qty{364}{\celsius}} más caliente que un
basalto primario de las dorsales actuales.
\end{solution}
